\section{Regularity of initial birth regions}
\label{sec:area_law_initial_region_regularity}

Retain the actual initial family of \cite[Section~11]{OpenAI2026AreaLaw},
with arbitrary origin, finite endpoint set and valid residues,
$C\ge 2$ and $k_0\ge 50{,}000{,}000$.

% Source: 10-geometry.tex, prop:two-families, lines 154--177,
% 212--218 and 299--323, especially 308--316;
% revision adc7f1241b42e322a6451854ab7e4b4c146bf78a.
% This auxiliary assertion concerns the initial regions before repairs.

\begin{theorem}[Birth regions are closures of their interiors]
    \label{thm:area_law_initial_birth_regularity}
    \lean{TNLean.PEPS.AreaLaw.Geometry.initialBirthRegion_eq_closure_initialOpenRegion}
    \leanok
    \uses{def:area_law_initial_birth_regions,def:area_law_initial_open_regions}
    For every actual initial identifier $i$,
    \begin{align}
        X_i^0 & =\overline{\operatorname{int}(X_i^0)}.
    \end{align}
    The endpoint set may be empty. Empty regions and disconnected
    primary regions are included.
\end{theorem}
\begin{proof}\leanok
    \uses{thm:area_law_cell_fan_cover,
        thm:area_law_primary_fine_cell_cover,def:area_law_fan_runs,
        def:area_law_initial_birth_regions}
    Each actual fan triangle is a closed convex hull with nonzero
    determinant. Its two edge vectors are linearly independent and
    span the plane, so the three vertices affinely span the plane.
    Thus the triangle has nonempty interior. A convex set with
    nonempty interior has the same closure as its interior; closedness
    therefore gives the asserted identity for each triangle.

    One unsplit fan triangle is contained in each closed dyadic square.
    Its nonempty interior supplies nonempty interior of the square.
    The same convexity argument proves regularity of the square.

    A finite union of regular closed sets is regular closed. Indeed,
    each member's interior lies in the union's interior, so taking
    closures shows that the union lies in the closure of that interior.
    The reverse inclusion follows from closedness of the finite union.
    This argument also includes the empty union.

    The dummy is a finite union of closed coarse squares, each primary
    has its exact finite cover by closed nonbelt fine cells, and each
    run is a finite union of its actual fan triangles. The finite-union
    argument proves the identity in all three identifier cases.
\end{proof}
